\section{Notation, Problem Formulation and Setup}

Given an LLM $M$, an input query $\vec{s}$, and the model-generated response
$\vec{\hat{g}}$ of length $T$, our goal is twofold: (i) predict whether
$\vec{\hat{g}}$ is hallucinated, and (ii) localize which generated-token
regions show uncertainty signals in the hidden-state activations. We assume
a white-box setting -- i.e.\ full access to the internal states of $M$ --
and disallow external resources such as auxiliary LLMs, retrieval, or
repeated sampling, so that detection can run alongside generation in real
time.

\paragraph{Activation Tensors.} For an LLM $M$ with $L$ transformer layers
and hidden dimension $D$, we collect the post-residual hidden states across
all layers and all generated tokens, yielding an activation tensor
$A \in \mathbb{R}^{L \times T \times D}$, where $T$ is the response
length.\footnote{This definition can be readily extended to include the
input text as well.} We refer to the layer--token plane as the spatial axes
of $A$ and to the hidden dimension as its feature axis; see
Figure~\ref{fig:overview}. Throughout this paper, we use the hidden states
at each layer after the application of the feed-forward sub-layer and the
residual connection (see Appendix~\ref{app:extended} for details).

\paragraph{Token Windows.} For temporal localization, we additionally
consider local views of $A$ along the token axis. For a window size $W$ and
a position $t \in \{1, \dots, T\}$, the windowed activation tensor is
\[
    A^{(t)} \in \mathbb{R}^{L \times W \times D},
\]
obtained by slicing $A$ over $W$ consecutive tokens centered at $t$ (with
edge padding when $t$ is near a response boundary). The full response is
then covered by the sequence $\{A^{(t)}\}_{t=1}^{T}$ of overlapping windows,
each providing a local view of the activation patterns associated with a
short generated-token region.

\paragraph{Dataset construction.} Following \citet{orgad2024}, we consider
datasets $\mathcal{D} = \{(\vec{s}_i, \vec{g}_i)\}_{i=1}^m$ of
question--answer pairs, e.g., (``Who developed the theory of general
relativity?'', ``Albert Einstein''). For each query $\vec{s}_i$, the model
produces a response $\vec{\hat{g}}_i$ from which we extract the activation
tensor $A_i$. We assign a sequence-level binary hallucination label $y_i$
by comparing $\vec{\hat{g}}_i$ with $\vec{g}_i$, setting $y_i = 1$ when the
response is incorrect/hallucinated and $y_i = 0$ otherwise. The final
activation dataset is $\mathcal{A} = \{(A_i, y_i, M_i)\}_{i=1}^m$, where
$M_i$ identifies the LLM that produced $A_i$, and multiple distinct $M_i$'s
from a known collection $\mathcal{M}$ are allowed within the same dataset.
\emph{Crucially, no token- or window-level labels are provided}: supervision
is restricted to the sequence level.

\paragraph{Problem statement.} We seek a parametric model that, given an
activation tensor $A$, predicts (i) a sequence-level uncertainty score
$\hat{p} \in [0, 1]$ and (ii) per-window uncertainty scores
$\{\hat{p}^{(t)}\}_{t=1}^{T}$ that localize the generated-token regions
whose activations most contribute to the sequence-level prediction. The
sequence-level prediction is supervised directly by $y$, while the
per-window scores must be learned under \emph{weak supervision} from the
same sequence label. Beyond the in-domain setting, we additionally consider
the following more general and challenging scenarios: (i) Off-domain data
generalization: ATs are drawn from a known LLM $M \in \mathcal{M}$ but
correspond to tasks or domains unseen during training; (ii) Off-domain
model generalization: ATs are drawn from a previously unseen LLM
$M \notin \mathcal{M}$.
